\documentclass[aspectratio=169,10pt,english]{beamer}

\input{../../../_preambulo_beamer.tex}
\input{../../../_comandos_beamer.tex}

\selectlanguage{english}

\title{MA1001B - Statistical Modeling for Decision Making}
\subtitle{Lesson 14: Bernoulli Trials and the Binomial Model}
\author{}
\institute{Tecnol\'ogico de Monterrey}
\date{}

\begin{document}

\begin{frame}[plain]
  \vspace{1.2cm}
  {\Large\bfseries MA1001B - Statistical Modeling for Decision Making\par}
  \vspace{0.7cm}
  {\large Lesson 14: Bernoulli Trials and the Binomial Model\par}
  \vspace{0.7cm}
  {\color{TecRojo}\rule{\textwidth}{1.2pt}}
  \vspace{0.8cm}

  MA1001B Faculty\\[0.3cm]
  Term\\[0.3cm]
  Tecnol\'ogico de Monterrey
\end{frame}

\begin{frame}{From One Trial to a Count}
  \begin{alertblock}{Guiding question}
    If each of $n=20$ inspections is defective with chance $p=0.08$, what is
    $P(X=0)$, and what breaks if $p$ is not the same on every lot?
  \end{alertblock}

  \vspace{0.6em}
  By the end of this lesson, you can:
  \begin{itemize}
    \item state the Bernoulli trial and the binomial assumptions;
    \item compute $P(X=0)$ and $P(X\le 2)$;
    \item report $np$ and $np(1-p)$;
    \item explain why mixed lot types are not binomial.
  \end{itemize}

  \vfill
  \scriptsize All inspections used today are synthetic. No real company data are used.
\end{frame}

\begin{frame}{Bernoulli Trial, Then $n$ Independent Copies}
  \small
  One inspection: $P(\text{defective})=p=0.08$, $P(\text{clean})=q=0.92$.
  Repeat $n=20$ times, independently, with the same $p$. Then
  $X\sim\mathrm{Binomial}(n=20,p=0.08)$ and
  \[
    E[X]=np=1.6,\qquad \Var(X)=np(1-p)=1.472.
  \]
\end{frame}

\begin{frame}[fragile]{Step 1: Bernoulli Trials}
  \begin{columns}[T]
    \begin{column}{0.42\textwidth}
      \small
      \begin{lstlisting}
p = 0.08
q = 1 - p
n = 20
mean = n * p
      \end{lstlisting}

      \begin{block}{Verified output}
        $p=0.080000$\\
        $q=0.920000$\\
        $n=20$\\
        $np=1.600000$
      \end{block}
    \end{column}
    \begin{column}{0.58\textwidth}
      \centering
      \includegraphics[width=\textwidth]{../figures/binomial_01_bernoulli_trials.png}
      \vspace{0.2em}
      \scriptsize One Bernoulli inspection from Step 1
    \end{column}
  \end{columns}
\end{frame}

\begin{frame}{Binomial Left Tail}
  \small
  \[
    P(X=0)=(0.92)^{20}=0.188693,
  \]
  \[
    P(X\le 2)=0.787946.
  \]
  These values use \texttt{scipy.stats.binom} with $n=20$ and $p=0.08$.
\end{frame}

\begin{frame}[fragile]{Step 2: PMF, CDF, Moments}
  \begin{columns}[T]
    \begin{column}{0.40\textwidth}
      \small
      \begin{lstlisting}
model = stats.binom(n=20, p=0.08)
p0 = model.pmf(0)
p2 = model.cdf(2)
      \end{lstlisting}

      \begin{block}{Verified output}
        $P(X=0)=0.188693$\\
        $P(X\le 2)=0.787946$\\
        $\Var(X)=1.472000$
      \end{block}
    \end{column}
    \begin{column}{0.60\textwidth}
      \centering
      \includegraphics[width=\textwidth]{../figures/binomial_02_pmf_cdf.png}
      \vspace{0.2em}
      \scriptsize Binomial PMF from Step 2
    \end{column}
  \end{columns}
\end{frame}

\begin{frame}{Step 3: $p$ Not Constant Across Lots}
  \begin{columns}[T]
    \begin{column}{0.42\textwidth}
      \small
      Half the lots have $p=0.02$ and half have $p=0.14$. The average $p$ is
      still $0.08$, and $E[X]$ is still $1.600000$.

      \vspace{0.4em}
      Mixture $P(X=0)=0.358291$ versus binomial $0.188693$. Mixture variance
      $2.840000$ versus $1.472000$.

      \vspace{0.3em}
      \textbf{Limit:} $p$ not constant across lots breaks the binomial.
      Lesson 15 compares two binomial policies with the same mean.
    \end{column}
    \begin{column}{0.58\textwidth}
      \centering
      \includegraphics[width=\textwidth]{../figures/binomial_03_nonconstant_p.png}
      \vspace{0.2em}
      \scriptsize Mixture versus binomial from Step 3
    \end{column}
  \end{columns}
\end{frame}

\begin{frame}{Concept Check}
  \begin{enumerate}
    \item What are the two outcomes of one Bernoulli inspection here?
    \item Compute $np$ for $n=20$ and $p=0.08$.
    \item Why is $P(X=0)=(0.92)^{20}$?
    \item Why does averaging $p=0.02$ and $p=0.14$ fail to restore the binomial?
  \end{enumerate}

  \vspace{0.6em}
  \begin{block}{Connection to Lesson 15}
    The session can continue directly into two binomial policies that share
    $np$ but not the same tail risk.
  \end{block}

  \vfill
  \scriptsize Source: OpenStax, \textit{Introductory Business Statistics 2e},
  Section 4.2. CC BY-NC-SA.
\end{frame}

\appendix



\end{document}
